\documentclass[11pt]{article}  %% nolint
\usepackage[margin=1in]{geometry}  %% nolint
\usepackage{amsmath,amssymb,booktabs,hyperref}
% ASLA paper macros --- GENERATED FILE, DO NOT EDIT BY HAND.
%
% Every value is read from a committed results JSON by key path via
% scripts/render_macros.py. Editing a number here breaks the provenance
% chain that exists because a retyped 4.22 survived into five documents
% and a draft before it was found to be wrong.
%
% Regenerate with:  python scripts/render_macros.py
%
% generated:      2026-09-23T07:02:30Z
% source commit:  5f4e07fae65f77657307cce966cdb97c7aa2f5c5-dirty

\newcommand{\residRatio}{2.90}  % theory_v2.json
\newcommand{\varInflation}{8.40}  % theory_v2.json
\newcommand{\residRatioSuperseded}{4.22}  % theory_v2.json (WRONG, kept for the correction note)
\newcommand{\varInflationSuperseded}{17.8}  % theory_v2.json (WRONG, kept for the correction note)
\newcommand{\singleRate}{1.33}  % a1_independent_rederivation.json
\newcommand{\projRate}{5.33}  % a1_independent_rederivation.json
\newcommand{\maxExcessWinning}{0.0}  % a2_dose_response.json
\newcommand{\nMatchedLadders}{48}  % a2_dose_response.json
\newcommand{\leverArmRho}{+0.564}  % a2_dose_response.json
\newcommand{\leverArmP}{3.2e-28}  % a2_dose_response.json
\newcommand{\budgetCountRho}{-0.136}  % a2_dose_response.json
\newcommand{\allocExcessMin}{+0.028}  % allocation_1B.json
\newcommand{\allocExcessMax}{+0.117}  % allocation_1B.json
\newcommand{\maxShareGap}{0.00058}  % allocation_1B.json
\newcommand{\bootCoverage}{0.00}  % abstention_1B.json
\newcommand{\confCoverage}{0.88}  % abstention_1B.json
\newcommand{\bootErrorToWidth}{4.73}  % abstention_1B.json
\newcommand{\centringInSeedSE}{119}  % abstention_1B.json
\newcommand{\bootEffectiveDelta}{0.500}  % abstention_1B.json
\newcommand{\confEffectiveDelta}{0.060}  % abstention_1B.json
\newcommand{\certifiedPairs}{292}  % abstention_1B.json
\newcommand{\totalPairs}{300}  % abstention_1B.json
\newcommand{\certifiedErrors}{0}  % abstention_1B.json
\newcommand{\correctCallsAbstained}{1.7}  % abstention_1B.json
\newcommand{\unresolvedPairs}{437}  % resolution_study.json
\newcommand{\adjacentPairs}{468}  % resolution_study.json
\newcommand{\unresolvedPct}{93.4}  % resolution_study.json
\newcommand{\unidentifiablePairs}{28}  % resolution_study.json
\newcommand{\overshoot}{+0.0877}  % curvature_gate.json
\newcommand{\curvatureRho}{-0.887}  % curvature_gate.json
\newcommand{\curvatureReduction}{22}  % curvature_gate.json
\newcommand{\tokensPerParamReduction}{60}  % curvature_gate.json
\newcommand{\targetTokensPerParam}{85.0}  % curvature_gate.json
\newcommand{\tokensPerParamCorr}{-0.80}  % curvature_gate.json
\newcommand{\confoundSpearman}{0.9977}  % curvature_gate.json
\newcommand{\determinedFracBonf}{82}  % target_scoring.json
\newcommand{\determinedFracBH}{98}  % target_scoring.json
\newcommand{\excessObserved}{2.33}  % target_scoring.json
\newcommand{\excessDetermined}{0.00}  % target_scoring.json
\newcommand{\excessExpected}{1.82}  % target_scoring.json
\newcommand{\welchDfMedian}{2.90}  % target_scoring.json
\newcommand{\welchDfMin}{2.00}  % target_scoring.json
\newcommand{\welchDfMax}{4.00}  % target_scoring.json
\newcommand{\tRatioAssumed}{3.63}  % target_scoring.json (superseded: assumed 4 df)
\newcommand{\tRatioMedian}{6.77}  % target_scoring.json
\newcommand{\tRatioPTen}{3.69}  % target_scoring.json
\newcommand{\tRatioPNinety}{18.27}  % target_scoring.json
\newcommand{\knownAnswerPerSeed}{0.8033}  % datadecide_known_answer.json
\newcommand{\knownAnswerSeedMean}{0.83}  % datadecide_known_answer.json
\newcommand{\knownAnswerPublished}{0.80}  % datadecide_known_answer.json
\newcommand{\focalScale}{150M}  % datadecide_known_answer.json
\newcommand{\targetScaleLabel}{1B}  % datadecide_known_answer.json
\newcommand{\maxDfAtThreeSeeds}{4}  % target_scoring.json
\newcommand{\expectedGap}{1.82}  % expected_error.json
\newcommand{\expectedCandLow}{-0.91}  % expected_error.json
\newcommand{\expectedCandHigh}{+4.55}  % expected_error.json
\newcommand{\expectedCandP}{0.19}  % expected_error.json
\newcommand{\expectedPairLow}{+0.27}  % expected_error.json
\newcommand{\expectedPairHigh}{+3.58}  % expected_error.json
\newcommand{\expectedPairP}{0.024}  % expected_error.json
\newcommand{\dependenceFactor}{1.63}  % expected_error.json
\newcommand{\pubObserved}{0.830}  % expected_error.json
\newcommand{\pubPosterior}{0.811}  % expected_error.json
\newcommand{\pubDisattenuated}{0.867}  % expected_error.json
\newcommand{\pubReliability}{0.950}  % expected_error.json
\newcommand{\pubDetFrac}{4}  % expected_error.json
\newcommand{\relSmallGaps}{0.860}  % expected_error.json
\newcommand{\relLargeGaps}{0.998}  % expected_error.json
\newcommand{\doseRhoObserved}{+0.860}  % expected_error.json
\newcommand{\doseRhoExpected}{+0.869}  % expected_error.json
\newcommand{\doseDesigns}{55}  % expected_error.json
\newcommand{\shippedProjFlips}{16}  % target_scoring.json (superseded; does not reproduce)
\newcommand{\recomputedProjFlips}{10}  % target_scoring.json
\newcommand{\recomputedSingleFlips}{3}  % target_scoring.json
\newcommand{\equivBoundPP}{1.22}  % robustness.json
\newcommand{\publishedOverall}{0.830}  % robustness.json
\newcommand{\publishedDeterminedBonf}{0.917}  % robustness.json
\newcommand{\publishedShiftBonf}{8.7}  % robustness.json
\newcommand{\publishedDeterminedBH}{0.861}  % robustness.json
\newcommand{\publishedShiftBH}{3.1}  % robustness.json
\newcommand{\publishedDetFracBonf}{4}  % robustness.json
\newcommand{\publishedDetFracBH}{62}  % robustness.json
\newcommand{\varDiagnosisRatio}{1.55}  % theory_check.json
\newcommand{\gapSdUnderstated}{1.25}  % theory_check.json
\newcommand{\sigmaJensenFactor}{2.06}  % theory_check.json
\newcommand{\nCandidates}{25}  % expected_error.json
\newcommand{\expectedSE}{1.39}  % expected_error.json
\newcommand{\bootCandLow}{-0.90}  % expected_error.json
\newcommand{\bootCandHigh}{+5.82}  % expected_error.json
\newcommand{\bootCandP}{0.21}  % expected_error.json
\newcommand{\uniqueSetP}{0.14}  % expected_error.json
\newcommand{\zetaRatio}{28}  % expected_error.json
\newcommand{\calibU}{1.03}  % expected_error.json
\newcommand{\calibBoot}{1.28}  % expected_error.json
\newcommand{\calibJack}{1.20}  % expected_error.json
\newcommand{\calibUnique}{0.99}  % expected_error.json
\newcommand{\powerNeeded}{87}  % expected_error.json
\newcommand{\powerNeededOne}{260}  % expected_error.json
\newcommand{\powerNeededTwo}{73}  % expected_error.json
\newcommand{\powerNeededFive}{18}  % expected_error.json
\newcommand{\powerAtObserved}{26}  % expected_error.json
\newcommand{\rescoreRowsLow}{Checkpoint-augmented & $+1.67$ & $+1.44$ & $[-0.26,\ +3.15]$ & weakens \\ Empirical-Bayes shrinkage & $+2.00$ & $+1.49$ & $[-1.30,\ +4.27]$ & weakens \\ Ensemble & $+9.00$ & $+8.53$ & $[+3.00,\ +14.06]$ & survives \\ Projection & $+2.33$ & $+1.82$ & $[-0.91,\ +4.55]$ & weakens \\ Shared exponent & $+0.00$ & $-0.23$ & $[-0.93,\ +0.47]$ & no observed difference \\}  % rescoring.json
\newcommand{\rescoreRowsHigh}{Checkpoint-augmented & $-0.33$ & $-0.01$ & $[-2.58,\ +2.56]$ & vanishes \\ Empirical-Bayes shrinkage & $+0.33$ & $+0.37$ & $[-2.03,\ +2.76]$ & weakens \\ Ensemble & $-1.33$ & $-0.54$ & $[-2.78,\ +1.69]$ & weakens \\ Projection & $+0.33$ & $+0.37$ & $[-2.03,\ +2.76]$ & weakens \\ Shared exponent & $+1.00$ & $+1.00$ & $[-2.17,\ +4.17]$ & weakens \\}  % rescoring.json
\newcommand{\rescoreDetLow}{82}  % rescoring.json
\newcommand{\rescoreDetHigh}{3}  % rescoring.json
\newcommand{\rescoreRowsAllocation}{0.01 of budget & $+10.00$ & $+9.72$ & $[+3.71,\ +15.72]$ & survives \\ 0.03 of budget & $+9.67$ & $+9.50$ & $[+2.33,\ +16.67]$ & survives \\ 0.1 of budget & $+8.00$ & $+7.84$ & $[+0.84,\ +14.84]$ & survives \\ 0.3 of budget & $+4.33$ & $+4.18$ & $[-0.02,\ +8.38]$ & weakens \\ 0.5 of budget & $+4.33$ & $+4.18$ & $[-0.02,\ +8.38]$ & weakens \\ 1 of budget & $+2.00$ & $+1.96$ & $[-0.37,\ +4.30]$ & weakens \\}  % rescoring.json
\newcommand{\certNormalN}{292}  % rescoring.json
\newcommand{\certNormalExpected}{0.26}  % rescoring.json
\newcommand{\certStudentTN}{280}  % rescoring.json
\newcommand{\certStudentTExpected}{0.01}  % rescoring.json
\newcommand{\certNormalMinP}{0.83}  % rescoring.json
\newcommand{\gateExact}{30}  % published_comparisons.json
\newcommand{\gateRows}{88}  % published_comparisons.json
\newcommand{\gateMaxError}{0.88}  % published_comparisons.json
\newcommand{\dOneVariants}{8}  % published_comparisons.json
\newcommand{\dOneMacroExcl}{1}  % published_comparisons.json
\newcommand{\dOneWorstDiff}{-15.2}  % published_comparisons.json
\newcommand{\dOneWorstLow}{-28.1}  % published_comparisons.json
\newcommand{\dOneWorstHigh}{-2.3}  % published_comparisons.json
\newcommand{\dOneTypicalHalfWidth}{8}  % published_comparisons.json
\newcommand{\dOneAll}{176}  % published_comparisons.json
\newcommand{\dOneAllBetter}{25}  % published_comparisons.json
\newcommand{\dOneAllWorse}{28}  % published_comparisons.json
\newcommand{\dTwoScales}{13}  % published_comparisons.json
\newcommand{\dTwoMacroBetter}{1}  % published_comparisons.json
\newcommand{\dTwoMacroWorse}{1}  % published_comparisons.json
\newcommand{\dTwoMacroPointWorse}{5}  % published_comparisons.json
\newcommand{\dTwoAll}{143}  % published_comparisons.json
\newcommand{\dTwoAllBetter}{56}  % published_comparisons.json
\newcommand{\dTwoAllWorse}{9}  % published_comparisons.json
\newcommand{\curveComparisons}{495}  % published_comparisons.json
\newcommand{\curveZetaMedian}{0.068}  % published_comparisons.json
\newcommand{\curveMedTen}{1.45}  % published_comparisons.json
\newcommand{\curveMedTwentyFive}{2.03}  % published_comparisons.json
\newcommand{\curveMedHundred}{3.79}  % published_comparisons.json
\newcommand{\curveLowTwentyFive}{1.33}  % published_comparisons.json
\newcommand{\curveHighTwentyFive}{2.55}  % published_comparisons.json
\newcommand{\curveCFour}{1.62}  % published_comparisons.json
\newcommand{\curveCheckWorstLarge}{11}  % published_comparisons.json
\newcommand{\curveCheckSmallest}{44}  % published_comparisons.json
\newcommand{\curveCheckSmallestM}{8}  % published_comparisons.json
\newcommand{\pInvariantMin}{0.14}  % expected_error.json
\newcommand{\pInvariantMax}{0.27}  % expected_error.json
\newcommand{\floorN}{25}  % floor_profile.json
\newcommand{\floorCFourExcl}{25}  % floor_profile.json
\newcommand{\floorErrExcl}{0}  % floor_profile.json
\newcommand{\floorErrWidthRel}{90}  % floor_profile.json
\newcommand{\floorCFourWidthRel}{17}  % floor_profile.json
\newcommand{\floorCondRatio}{184}  % floor_profile.json
\newcommand{\floorMisspecCFour}{8.4}  % floor_profile.json
\newcommand{\floorNewWidth}{0.75}  % floor_profile.json
\newcommand{\floorOldWidth}{0.44}  % floor_identifiability.json
\newcommand{\floorFCrit}{5.32}  % floor_profile.json
\newcommand{\thmDesigns}{2000}  % fit_structure.json
\newcommand{\thmCorr}{0.999997}  % fit_structure.json
\newcommand{\thmExcessThree}{0.0024}  % fit_structure.json
\newcommand{\thmSolverExcess}{1.8}  % fit_structure.json
\newcommand{\couplingPairs}{66}  % seed_coupling.json
\newcommand{\couplingMedianCFour}{-0.004}  % seed_coupling.json
\newcommand{\couplingMedianErr}{+0.008}  % seed_coupling.json
\newcommand{\couplingRejectCFour}{4.5}  % seed_coupling.json
\newcommand{\couplingRejectErr}{0.0}  % seed_coupling.json
\newcommand{\couplingFocusRho}{0.30}  % seed_coupling.json
\newcommand{\couplingFocusP}{0.13}  % seed_coupling.json
\newcommand{\polyCells}{10}  % moderated_variance.json
\newcommand{\polyBetter}{10}  % moderated_variance.json
\newcommand{\polyRatioRaw}{0.890}  % moderated_variance.json
\newcommand{\polyRatioMod}{0.926}  % moderated_variance.json
\newcommand{\polyRmseRaw}{0.72}  % moderated_variance.json
\newcommand{\polyRmseMod}{0.50}  % moderated_variance.json
\newcommand{\nullTests}{1000000}  % moderated_variance.json
\newcommand{\nullModFive}{1.22}  % moderated_variance.json
\newcommand{\nullModBonf}{4.9}  % moderated_variance.json
\newcommand{\nullRawFive}{0.69}  % moderated_variance.json
\newcommand{\nullRawBonf}{0.44}  % moderated_variance.json
\newcommand{\brierRaw}{0.0366}  % moderated_variance.json
\newcommand{\brierModT}{0.0363}  % moderated_variance.json
\newcommand{\brierAdopted}{0.0362}  % moderated_variance.json
\newcommand{\detModCFour}{96}  % moderated_variance.json
\newcommand{\detModErr}{46}  % moderated_variance.json
\newcommand{\detRawErrFiveThirty}{3}  % moderated_variance.json
\newcommand{\calibrationRows}{Projection & $+2.35$ & $-0.26$ & $-0.44$ & 0.86 & 0.88 & 0.90 & 0.98 & 0.68 & 0.03 \\ Shared exponent & $+0.68$ & $-0.42$ & $-0.08$ & 0.90 & 0.85 & 0.93 & 0.97 & 0.13 & 0.01 \\ EB shrinkage & $+1.64$ & $-0.62$ & $-0.43$ & 0.89 & 0.89 & 0.93 & 0.99 & 0.45 & 0.03 \\ Ensemble & $+8.05$ & $+3.11$ & $-0.36$ & 0.89 & 0.89 & 0.92 & 0.95 & 0.99 & 0.35 \\ Checkpoint-aug. & $+1.68$ & $+6.04$ & $-0.11$ & 0.91 & 0.87 & 0.98 & 0.92 & 0.73 & 0.87 \\ Top rung 150M (planted) & $-0.44$ & $+3.64$ & $-0.02$ & 0.91 & 0.85 & 0.96 & 0.94 & 0.15 & 0.80 \\ Top rung 90M (planted) & $+0.18$ & $+5.28$ & $-0.10$ & 0.92 & 0.85 & 0.96 & 0.95 & 0.05 & 0.94 \\ Top rung 60M (planted) & $+2.05$ & $+8.02$ & $-0.09$ & 0.95 & 0.86 & 0.95 & 0.92 & 0.65 & 0.99 \\ Top rung 20M (planted) & $+7.33$ & $+14.43$ & $-0.04$ & 0.96 & 0.90 & 0.99 & 0.93 & 1.00 & 1.00 \\}  % calibration.json
\newcommand{\calFixedWorlds}{150}  % calibration.json
\newcommand{\calRandomWorlds}{300}  % calibration.json
\newcommand{\calBootDraws}{50}  % calibration.json
\newcommand{\calFixedK}{2.26}  % calibration.json
\newcommand{\calJackK}{1.73}  % calibration.json
\newcommand{\calUK}{3.09}  % calibration.json
\newcommand{\calJackHalf}{3.18}  % calibration.json
\newcommand{\calUHalf}{4.10}  % calibration.json
\newcommand{\calFixedCovMin}{0.90}  % calibration.json
\newcommand{\calFixedCovMax}{0.99}  % calibration.json
\newcommand{\calJackCovMin}{0.92}  % calibration.json
\newcommand{\calJackCovMax}{0.99}  % calibration.json
\newcommand{\calUCovMin}{0.85}  % calibration.json
\newcommand{\calUCovMax}{0.90}  % calibration.json
\newcommand{\calProjFixedCov}{0.90}  % calibration.json
\newcommand{\calMaxBias}{0.44}  % calibration.json
\newcommand{\calCouplingShift}{0.12}  % calibration.json
\newcommand{\headEstimate}{1.82}  % calibration.json
\newcommand{\headFixedLow}{+0.15}  % calibration.json
\newcommand{\headFixedHigh}{+3.49}  % calibration.json
\newcommand{\headFixedP}{0.030}  % calibration.json
\newcommand{\headRandomLow}{-1.02}  % calibration.json
\newcommand{\headRandomHigh}{+4.65}  % calibration.json
\newcommand{\headRandomP}{0.18}  % calibration.json
\newcommand{\ensFixedP}{0.0007}  % calibration.json
\newcommand{\ensRandomP}{0.012}  % calibration.json
\newcommand{\restatedPrimaryRows}{projection vs single-scale, c4\_en\_bits\_per\_token & $+1.82$ & $[-1.0,\ +4.7]$ & 0.18 & --- \\ shared\_exponent vs single-scale, c4\_en\_bits\_per\_token & $-0.23$ & $[-1.1,\ +0.6]$ & 0.59 & --- \\ eb\_shrinkage vs single-scale, c4\_en\_bits\_per\_token & $+1.49$ & $[-1.3,\ +4.3]$ & 0.28 & --- \\ ensemble vs single-scale, c4\_en\_bits\_per\_token & $+8.53$ & $[+2.8,\ +14.3]$ & 0.012 & nominal \\ checkpoint\_augmented vs single-scale, c4\_en\_bits\_per\_token & $+1.44$ & $[-0.9,\ +3.8]$ & 0.22 & --- \\ projection vs single-scale, olmes\_macro\_error & $+0.37$ & $[-2.3,\ +3.0]$ & 0.78 & --- \\ shared\_exponent vs single-scale, olmes\_macro\_error & $+1.00$ & $[-2.3,\ +4.3]$ & 0.53 & --- \\ eb\_shrinkage vs single-scale, olmes\_macro\_error & $+0.37$ & $[-2.3,\ +3.0]$ & 0.78 & --- \\ ensemble vs single-scale, olmes\_macro\_error & $-0.54$ & $[-2.9,\ +1.8]$ & 0.64 & --- \\ checkpoint\_augmented vs single-scale, olmes\_macro\_error & $-0.01$ & $[-2.9,\ +2.8]$ & 0.99 & --- \\ optimised allocation at 0.01 of budget vs single-scale, c4 & $+9.72$ & $[+3.7,\ +15.8]$ & 0.0077 & nominal \\ optimised allocation at 0.03 of budget vs single-scale, c4 & $+9.50$ & $[+2.5,\ +16.5]$ & 0.019 & nominal \\ optimised allocation at 0.1 of budget vs single-scale, c4 & $+7.84$ & $[+1.1,\ +14.6]$ & 0.029 & nominal \\ optimised allocation at 0.3 of budget vs single-scale, c4 & $+4.18$ & $[-0.1,\ +8.5]$ & 0.057 & --- \\ optimised allocation at 0.5 of budget vs single-scale, c4 & $+4.18$ & $[-0.1,\ +8.5]$ & 0.057 & --- \\ optimised allocation at 1 of budget vs single-scale, c4 & $+1.96$ & $[-0.7,\ +4.6]$ & 0.13 & --- \\ projection excess rises with log lever arm (expected-error scoring) & --- & --- & 8.2e-20 & Holm \\ DataDecide 2\_param-default vs single-scale at 750M (exploratory: open discrepancy) & $-8.22$ & $[-19.5,\ +3.1]$ & 0.13 & --- \\ DataDecide 3\_param-1\_step vs single-scale at 750M (exploratory: open discrepancy) & $+1.11$ & $[-6.4,\ +8.7]$ & 0.77 & --- \\ DataDecide 3\_param-default-helper\_points-step2=0.5 vs single-scale at 750M (exploratory: open discrepancy) & $-1.56$ & $[-10.1,\ +7.0]$ & 0.71 & --- \\ DataDecide 3\_param-default-helper\_points vs single-scale at 750M (exploratory: open discrepancy) & $-0.89$ & $[-8.2,\ +6.5]$ & 0.81 & --- \\ DataDecide 3\_param-default-step2=0.5 vs single-scale at 750M (exploratory: open discrepancy) & $-1.89$ & $[-10.4,\ +6.6]$ & 0.65 & --- \\ DataDecide 3\_param-default vs single-scale at 750M (exploratory: open discrepancy) & $-1.89$ & $[-10.6,\ +6.8]$ & 0.66 & --- \\ DataDecide 5\_param-1\_step-ai2 vs single-scale at 750M (exploratory: open discrepancy) & $-15.22$ & $[-27.8,\ -2.7]$ & 0.027 & nominal \\ DataDecide 5\_param-ai2 vs single-scale at 750M (exploratory: open discrepancy) & $-2.22$ & $[-11.3,\ +6.8]$ & 0.61 & --- \\ Correct Prob vs accuracy at 10M, OLMES macro & $-5.44$ & $[-17.8,\ +6.9]$ & 0.36 & --- \\ Correct Prob vs accuracy at 14M, OLMES macro & $+2.33$ & $[-8.8,\ +13.4]$ & 0.67 & --- \\ Correct Prob vs accuracy at 150M, OLMES macro & $-1.67$ & $[-5.5,\ +2.2]$ & 0.37 & --- \\ Correct Prob vs accuracy at 16M, OLMES macro & $+6.44$ & $[-3.4,\ +16.2]$ & 0.17 & --- \\ Correct Prob vs accuracy at 20M, OLMES macro & $+4.44$ & $[-5.9,\ +14.7]$ & 0.37 & --- \\ Correct Prob vs accuracy at 300M, OLMES macro & $+1.00$ & $[-4.7,\ +6.7]$ & 0.72 & --- \\ Correct Prob vs accuracy at 4M, OLMES macro & $+29.11$ & $[+17.2,\ +41.0]$ & 0.0021 & nominal \\ Correct Prob vs accuracy at 530M, OLMES macro & $+0.56$ & $[-4.5,\ +5.6]$ & 0.83 & --- \\ Correct Prob vs accuracy at 60M, OLMES macro & $+4.00$ & $[-3.0,\ +11.0]$ & 0.24 & --- \\ Correct Prob vs accuracy at 6M, OLMES macro & $-5.11$ & $[-16.2,\ +5.9]$ & 0.34 & --- \\ Correct Prob vs accuracy at 750M, OLMES macro & $+1.22$ & $[-3.9,\ +6.4]$ & 0.63 & --- \\ Correct Prob vs accuracy at 8M, OLMES macro & $-13.78$ & $[-25.5,\ -2.0]$ & 0.029 & nominal \\ Correct Prob vs accuracy at 90M, OLMES macro & $-1.56$ & $[-5.7,\ +2.6]$ & 0.44 & --- \\}  % restated_claims.json
\newcommand{\restatedPrimaryM}{38}  % restated_claims.json
\newcommand{\restatedPrimaryNominal}{8}  % restated_claims.json
\newcommand{\restatedPrimaryHolm}{1}  % restated_claims.json
\newcommand{\restatedPrimaryBY}{1}  % restated_claims.json
\newcommand{\restatedFixedRows}{projection vs single-scale, c4\_en\_bits\_per\_token (these 25 recipes) & $+1.82$ & $[+0.1,\ +3.6]$ & 0.038 & nominal \\ shared\_exponent vs single-scale, c4\_en\_bits\_per\_token (these 25 recipes) & $-0.23$ & $[-1.7,\ +1.2]$ & 0.74 & --- \\ eb\_shrinkage vs single-scale, c4\_en\_bits\_per\_token (these 25 recipes) & $+1.49$ & $[+0.2,\ +2.7]$ & 0.019 & nominal \\ ensemble vs single-scale, c4\_en\_bits\_per\_token (these 25 recipes) & $+8.53$ & $[+5.1,\ +11.9]$ & 0.00074 & Holm \\ checkpoint\_augmented vs single-scale, c4\_en\_bits\_per\_token (these 25 recipes) & $+1.44$ & $[+0.4,\ +2.5]$ & 0.0076 & nominal \\ projection vs single-scale, olmes\_macro\_error (these 25 recipes) & $+0.37$ & $[-2.1,\ +2.8]$ & 0.75 & --- \\ shared\_exponent vs single-scale, olmes\_macro\_error (these 25 recipes) & $+1.00$ & $[-1.5,\ +3.5]$ & 0.41 & --- \\ eb\_shrinkage vs single-scale, olmes\_macro\_error (these 25 recipes) & $+0.37$ & $[-2.1,\ +2.8]$ & 0.75 & --- \\ ensemble vs single-scale, olmes\_macro\_error (these 25 recipes) & $-0.54$ & $[-3.9,\ +2.8]$ & 0.73 & --- \\ checkpoint\_augmented vs single-scale, olmes\_macro\_error (these 25 recipes) & $-0.01$ & $[-3.1,\ +3.1]$ & 1 & --- \\}  % restated_claims.json
\newcommand{\restatedFixedM}{10}  % restated_claims.json
\newcommand{\restatedFixedNominal}{4}  % restated_claims.json
\newcommand{\restatedFixedHolm}{1}  % restated_claims.json
\newcommand{\restatedFixedBY}{1}  % restated_claims.json
\newcommand{\doseJackP}{8.2e-20}  % expected_error.json
\newcommand{\doseJackLow}{+0.78}  % expected_error.json
\newcommand{\doseJackHigh}{+0.92}  % expected_error.json


\title{The Reference Is Part of the Measurement:\\
What Decision-Accuracy Comparisons Can Resolve at Realistic Candidate Counts}
\author{}
\date{}

\begin{document}
\maketitle

\begin{abstract}
Small-scale experiments are routinely used to choose between training
interventions, and selection rules are compared by \emph{decision accuracy}:
agreement with an ordering measured at a larger target scale. We give a
calibrated inference procedure for such comparisons and ask what it can
establish at realistic candidate counts. The procedure has three parts. The
target reference's uncertainty is propagated with expected-error scoring, using
variances moderated by empirical Bayes and validated against ten-seed ground
truth. The unit of inference is the candidate, not the pair. And two estimands
are separated: a fixed-candidate estimand for these recipes and a
random-candidate estimand for recipes in general. Critical values come from an
end-to-end semi-synthetic benchmark built from DataDecide. Validation changed
the procedure twice. The textbook moderated $t$ test is anti-conservative on
real few-seed data, by a factor of \nullModBonf\ at a Bonferroni level. Wald
intervals from unbiased candidate-level standard errors undercover, because
those standard errors are noisy at \nCandidates\ candidates. Under the calibrated
procedure, projection's excess mis-selection over single-scale ranking on C4
(\headEstimate\ points) is significant for these recipes ($p=\headFixedP$) but
not as a claim about recipes in general ($p=\headRandomP$). Across the paper's
\restatedPrimaryM\ inferential claims, \restatedPrimaryNominal\ are nominally
significant and \restatedPrimaryBY\ survives correction for dependent tests: the
growth of projection error with the extrapolation lever arm. A closed-form
correction converts pair-resampled intervals to candidate-level ones.
\end{abstract}

\paragraph{Contributions.} (i) A calibrated procedure for comparing selection
rules: moderated reference variances, candidate-level inference, and two
explicit estimands, each validated against known answers. (ii) Two negative
results about standard tools: the moderated $t$ test is anti-conservative on
real few-seed data, and unbiased candidate-level standard errors give
undercovering Wald intervals at realistic candidate counts. (iii) A power
calculation: a two-point difference needs about \powerNeededTwo\ candidates.
(iv) A closed-form correction $R(n)$ for intervals already computed by
resampling pairs. (v) Candidate-level re-tests of published comparisons, and an
account of which of this paper's own claims survive correction.

\section{How selection rules are compared}
\label{sec:setup}

A selection rule observes candidates at small compute and predicts their
ordering at a target budget. Its decision accuracy is the fraction of candidate
pairs whose predicted ordering matches a reference ordering measured at the
target. Writing $B(\cdot)$ for a benchmark score, the standard estimator is an
indicator on sign agreement between the small-scale and target gaps, averaged
over pairs.

The reference is a finite-seed estimate. For a pair whose target gap is small
relative to seed noise, the reference ordering is a coin flip, and every rule is
graded against that flip. The estimator does not distinguish such a pair from
one the reference resolves cleanly.

\section{Uncertainty-aware decision accuracy}
\label{sec:protocol}

For each pair we compute the target gap, its standard error from the seeds, and
a two-sided Welch $t$ test with Satterthwaite degrees of freedom computed per
pair. Two scoring rules follow.

\paragraph{Determined-subset scoring.} A pair is \emph{determined} if the test
rejects equality at a stated level with stated multiplicity control. Rules are
scored only on determined pairs. Assumption: cell means are approximately
Gaussian with unknown, possibly unequal variances.

\paragraph{Expected-error scoring.} Each pair is charged in expectation: a rule
predicting the observed order is charged $1-p$ and one predicting the opposite  %% nolint: algebra, not a measured value
is charged $p$, where $p$ is the probability the observed order is correct under
a Gaussian approximation to the difference in cell means. No prior beyond the
observed gap is used.

\paragraph{What neither rule fixes.} Both treat target cell means as unbiased
estimates of expected target performance. A systematic difference between them
--- a different evaluation draw, a different checkpoint --- survives both.

\paragraph{The noise model.} Three decisions determine every standard error in
this paper, and each was tested against a known answer before it was adopted.

\emph{Seeds are independent across scales.} DataDecide's seeds control
initialization and data order. If a seed label fixed the same order at every
size, single-scale ranking and the target would share noise, which would
flatter single-scale ranking. Across \couplingPairs\ pairs of scales, the
correlation between a seed's deviations at two scales has median
\couplingMedianCFour\ on C4 and \couplingMedianErr\ on OLMES error. Permutation
tests reject at \couplingRejectCFour\% and \couplingRejectErr\% of pairs, which
is what chance produces. The one pair that bears on the headline, 300M against  %% nolint: scale labels, not measured values
530M on C4, has $\rho=\couplingFocusRho$ ($p=\couplingFocusP$). Planting that  %% nolint: scale label, not a measured value
coupling in the benchmark below moves no estimate by more than
\calCouplingShift\ points.

\emph{Variances are moderated, but tail probabilities are not.} Three seeds give
a poor variance. Following Smyth (2004), we moderate each recipe's variance  %% nolint: citation year
toward a prior fitted across recipes at the same scale. We checked this on
PolyPythias, which has ten seeds per model. There, the tasks at each size and
step play the role of recipes, and three-seed estimates are compared with the
ten-seed truth. Moderation reduces the error in log standard deviation at
\polyBetter\ of \polyCells\ cells (median \polyRmseRaw\ to \polyRmseMod). It
only partly removes the low bias of a three-seed standard deviation, from
\polyRatioRaw\ to \polyRatioMod\ of the truth; we had predicted it would
remove it. The moderated $t$ test itself fails. Splitting ten seeds three
against three, so that the null is true by construction, it rejects at
\nullModFive\ times its nominal rate at $0.05$ and \nullModBonf\ times at the  %% nolint: stated test level
Bonferroni level for this design, over \nullTests\ tests. Raw Welch rejects at
\nullRawFive\ and \nullRawBonf\ times nominal. Target determination therefore
stays with the raw Welch test. Adopting the moderated $t$ test would have raised
the determined fraction on OLMES error from \detRawErrFiveThirty\% to
\detModErr\%, most of which would be false. The expected-error weight uses the
moderated standard error with the raw Welch degrees of freedom. Of three
candidate weights, that one scored the lowest Brier score against held-out
seeds (\brierAdopted, against \brierRaw\ for the unmoderated normal weight).

\section{Two estimands, and a calibrated test for each}
\label{sec:estimands}

A difference in decision accuracy between two rules answers one of two
questions, and they need different uncertainty.

\emph{Fixed-candidate: is the difference real for these recipes?} The recipes
are fixed; only seed noise is random. We resample every cell mean and
checkpoint parametrically from the moderated noise model around the observed
means, re-run every rule and the scoring on each draw, and take the spread.
Resampling three seeds instead would give about ten distinct multisets per cell.

\emph{Random-candidate: would the difference hold for new recipes?} The recipes
are a sample. The statistic is a mean over pairs of candidates, and its
uncertainty is candidate-level, as in \S\ref{sec:gap}.

\paragraph{End-to-end calibration.} To check both procedures we built a
semi-synthetic benchmark from DataDecide. It keeps DataDecide's candidates and
ladder, uses the observed seed-mean trajectories as the truth, and draws noise
from the moderated scale-dependent model, with AR(1) noise along each run's  %% nolint: model name AR(1)
checkpoints. Fixed-candidate worlds (\calFixedWorlds) hold the truth and
redraw the noise. Random-candidate worlds (\calRandomWorlds) draw a new set of
recipes by a smoothed bootstrap over the recipes' curves. In each world the true
difference in mis-selection is computed from the noiseless curves, and the full
procedure is run on the noisy data. Planted comparators, single-scale ranking at
lower rungs, supply known effects for power.

Table~\ref{tab:calibration} gives the results. Expected-error scoring is close
to unbiased for the true difference, with no bias larger than \calMaxBias\
points. The intervals are not calibrated at the textbook critical value. The
candidate-level standard errors are right on average but noisy at
\nCandidates\ candidates, so Wald intervals at 1.96 undercover: the U-statistic  %% nolint: the textbook critical value
covers \calUCovMin--\calUCovMax. We therefore calibrate the critical value on the
benchmark itself: \calFixedK\ for the fixed-candidate estimand and \calJackK\ for
the random-candidate one. For the random-candidate estimand the delete-one
jackknife gives narrower calibrated intervals than the U-statistic (median
half-width \calJackHalf\ against \calUHalf\ points), so we adopt it. For
shared-exponent and EB shrinkage, whose fits couple candidates, the jackknife
refits the whole rule. Calibration is imperfect by comparator. On held-out
worlds, calibrated coverage runs \calFixedCovMin--\calFixedCovMax\ for the
fixed-candidate estimand and \calJackCovMin--\calJackCovMax\ for the
random-candidate one. For projection, the fixed-candidate coverage is
\calProjFixedCov.

\begin{table}[t]
\centering
\scriptsize
\begin{tabular}{lrrrrrrrrr}
\toprule
& \multicolumn{2}{c}{true effect (pp)} & & \multicolumn{2}{c}{coverage at 1.96} & \multicolumn{2}{c}{calibrated, held out} & \multicolumn{2}{c}{calibrated power} \\  %% nolint: the textbook critical value
comparator & fixed & random & bias & fixed & random & fixed & random & fixed & random \\
\midrule
\calibrationRows
\bottomrule
\end{tabular}
\caption{End-to-end calibration on semi-synthetic DataDecide worlds (C4, fit to
300M, target 530M). Comparisons are against single-scale ranking at the top  %% nolint: scale labels
rung; positive means worse. Coverage at 1.96 uses the parametric bootstrap  %% nolint: the textbook critical value
(fixed) and the U-statistic (random); calibrated columns use the critical
values in the text, with the jackknife for the random-candidate estimand.
Planted comparators rank at a lower rung.}
\label{tab:calibration}
\end{table}

\section{Instrument validation}
\label{sec:instrument}

Before reporting any comparison we re-ran the DataDecide known-answer check at
the current commit: single-scale ranking at \focalScale\ predicting \targetScaleLabel\ reproduces
the published figure to within band (per-seed \knownAnswerPerSeed, seed-mean
\knownAnswerSeedMean, against a published \knownAnswerPublished).

\section{Two regimes, and where the split is not clean}
\label{sec:regimes}

The fraction of target pairs a reference resolves is a property of the metric,
not of the selection rule, and it separates sharply. On the low-noise metric (C4
bits per token), \determinedFracBonf\% of pairs are resolved under a
Bonferroni-corrected Welch test. On the accuracy metric behind the published
figure below, \pubDetFrac\% are. Decision-accuracy comparisons mean different
things in the two regimes.

The split is cleaner for the reference than for the fit. On the accuracy metric
the power-law floor is unidentified: its profile-likelihood interval includes
zero for every candidate and spans \floorErrWidthRel\% of the feasible range
(\S\ref{sec:fit}). We pre-registered that correcting the profile would
\emph{narrow} the C4 interval, and it widened it, from \floorOldWidth\ to
\floorNewWidth. The reason is that C4 residuals run \floorMisspecCFour\ times
the seed variance. The C4 floor still excludes zero for every candidate, but it
is poorly determined, and the misspecification that widens it is the same one
that dominates projection error. Low noise in the reference does not mean a
well-determined fit. The regimes classify how well the \emph{reference}
resolves decisions; they do not classify how well the extrapolation is
specified.

\section{Low-noise regime: what \nCandidates\ candidates can resolve}
\label{sec:gap}

On the primary design, projection's excess over single-scale ranking is
\excessObserved\ points under observed-target scoring and \expectedGap\ points
under expected-error scoring. The point estimate favours single-scale ranking,
and propagating reference uncertainty retains most of it, so the difference is
not an artifact of noisy labels.

Whether the difference is established is a separate question. Its unit is the
candidate: all pairs are formed from \nCandidates\ candidates, and each
candidate enters many of them. The statistic is a U-statistic of order two, and
its variance has a candidate-level and a pair-level component. The
candidate-level test gives a standard error of \expectedSE\ points and a 95\%  %% nolint: confidence level, a stated definition
interval of [\expectedCandLow, \expectedCandHigh] ($p=\expectedCandP$). As a
claim about recipes in general the comparison is underpowered, and no
difference is established. The calibrated procedure below reaches the same
conclusion for that question and a different one for these recipes.

For contrast, treating the pairs as independent gives
[\expectedPairLow, \expectedPairHigh] ($p=\expectedPairP$). This
anti-conservative version understates the standard error by a factor of
\dependenceFactor, and it would have reported a significant difference. The
factor is the size of the dependence correction. Any decision-accuracy
comparison that resamples or counts pairs as independent inherits it.

\paragraph{Under the calibrated procedure.} With the calibrated weights the
estimate is \headEstimate\ points. For these \nCandidates\ recipes the difference
is significant: the fixed-candidate interval is [\headFixedLow, \headFixedHigh]
($p=\headFixedP$). The caveat is that this estimand's held-out coverage for
projection is \calProjFixedCov, below nominal, so the interval is somewhat
optimistic. As a claim about recipes in general it is not established: the
random-candidate interval is [\headRandomLow, \headRandomHigh]
($p=\headRandomP$). Single-scale ranking is better than projection on these
recipes, and the data cannot say whether that transfers.

\paragraph{Why the estimator changed twice.} We pre-registered a candidate
bootstrap weighted by draw multiplicity. Simulation at this comparison's
variance components ($\zeta_2/\zeta_1=\zetaRatio$) showed it overstates the
standard error by a factor of \calibBoot. The unbiased U-statistic standard
error is within \calibU. The end-to-end benchmark then showed that even an
unbiased standard error gives undercovering intervals at this candidate count,
which led to the calibrated jackknife of \S\ref{sec:estimands}. Neither choice
was blind to the real-data results. The random-candidate conclusion does not
depend on it: every candidate-level estimator we computed gives $p$ between
\pInvariantMin\ and \pInvariantMax.

\paragraph{Power.} With the variance components held at their estimates, a
two-sided test at level 0.05 with power 0.8 needs \powerNeededTwo\ candidates to  %% nolint: stated test level and power
detect a true two-point difference, \powerNeededOne\ for one point and
\powerNeededFive\ for five. At \nCandidates\ candidates the power against the
observed difference is \powerAtObserved\%. This assumes a known standard error.
The calibrated procedure pays for estimating it, so these counts are lower
bounds. Published suites that compare selection rules typically have tens of
candidates. Differences of a few points in decision accuracy are below what
those suites can resolve.

\paragraph{On the determined subset.} Scoring only on resolved pairs gives an
excess of \excessDetermined\ points, but both rules make zero errors there, so
the comparison is empty. That subset rules out only differences above
\equivBoundPP\ points by the rule of three --- a bound set by subset size, not
by observed agreement. We report it as a diagnostic and not as evidence of
equivalence.

\section{High-noise regime: a published figure}
\label{sec:published}

DataDecide reports a pairwise decision accuracy of \knownAnswerPublished\ for
ranking at a single small scale; reproduced here it is \publishedOverall. Only
\pubDetFrac\% of its reference pairs are determined under Bonferroni. The
reference does not resolve most of the decisions the figure grades.

An earlier version of this work reported accuracy on the determined pairs as a
corrected figure. We withdraw it because it is selection-biased. Determined
pairs are the well-separated ones, where any rule scores near one, so the figure
measures an easier population rather than correcting the same quantity. We
report two estimators of the same quantity instead.

\emph{Posterior expected accuracy} charges each pair the probability that its
reference ordering is correct. It gives \pubPosterior, below the observed
figure: unresolved pairs contribute near one half and pull the average toward
chance. \emph{Disattenuated accuracy} inverts
$\text{observed}=ar+(1-a)(1-r)$ for the rule's accuracy $a$ at mean reference  %% nolint: mathematical constant
reliability $r=\pubReliability$. It gives \pubDisattenuated, above the observed
figure. It assumes the rule's errors are independent of the reference's, and
that fails here. Both err preferentially on small true gaps: mean reliability is
\relSmallGaps\ on the smallest-gap third of pairs and \relLargeGaps\ on the
largest.

The two estimators bracket the published figure and disagree in direction. That
disagreement is the finding. It is also a general warning. A standard
measurement-error correction applied to decision accuracy assumes the rule and
the reference err independently. Both err where true gaps are small, so the
correction moves the figure the wrong way exactly where it matters. The direction of the correction depends on an
assumption about how rule errors and reference errors co-occur, and the data do
not identify it. The robust statement is the determined fraction: at this
reference budget the published figure mostly grades decisions its reference
cannot make.

\section{Published comparisons at the candidate level}
\label{sec:published-comparisons}

Neither DataDecide nor \emph{Signal and Noise} reports an interval on a difference
between selection methods. DataDecide's shading is the standard deviation over
three small-model seeds for single-scale ranking; its scaling-law predictions
are single attempts. \emph{Signal and Noise} reports point estimates and
resamples only over final checkpoints. Neither resampled pairs or candidates.
The question is therefore not which unit their intervals used, but whether
their conclusions hold with candidates as the unit. We re-tested the
comparisons whose inputs were released.

\paragraph{Scaling-law variants against single-scale ranking (DataDecide).} The
released fits include per-recipe predictions and targets, and we pre-registered
reproducing the released decision accuracies exactly before testing. There is
an open discrepancy, which we are investigating with the authors. Using the
released target column, \gateExact\ of \gateRows\ task--variant rows match
our recomputation, and the largest difference is \gateMaxError. No convention
we tried (prediction column, target column or seed mean, tie rule) matches them
all. Until it is resolved, the re-test below is exploratory and makes no claim
about the paper's numbers. Against single-scale ranking at the largest
fitted size on the macro average, \dOneMacroExcl\ of \dOneVariants\ variants
is nominally different: a five-parameter one-step fit, which is worse by
\dOneWorstDiff\ points ([\dOneWorstLow, \dOneWorstHigh]). It does not survive
correction across the paper's claims (\S\ref{sec:survives}). No variant is
better, even nominally. The other intervals have a typical half-width of
\dOneTypicalHalfWidth\ points. Across all tasks, variants and comparators,
\dOneAllBetter\ of \dOneAll\ comparisons are nominally better and \dOneAllWorse\
nominally worse, uncorrected across comparisons. The published conclusion, that no scaling-law variant exceeds
single-scale ranking, is consistent with these data. Equivalence is not
established, and differences of several points in either direction are not
excluded.

\paragraph{Correct Prob against accuracy (DataDecide).} The paper reports that
ranking by the probability of the correct answer is at least as good as
accuracy ``for most small scales''. On the macro average, the candidate-level
test gives a nominal advantage at \dTwoMacroBetter\ of \dTwoScales\ small scales
and a nominal disadvantage at \dTwoMacroWorse, and neither survives correction.
The point estimate favours accuracy at \dTwoMacroPointWorse. Per task the claim has more support: of \dTwoAll\
task--scale cells, \dTwoAllBetter\ show a nominal advantage and
\dTwoAllWorse\ a nominal disadvantage, uncorrected across cells.

\paragraph{Uncheckable.} \emph{Signal and Noise}'s headline BPB comparison (the
30-task average) cannot be rebuilt: the paper names fifteen of the thirty tasks,  %% nolint: counts of named tasks in the source table
and the released data does not identify the rest. Its checkpoint-averaging
table does not state how source sizes are combined. Its high-SNR subsets and
DataDecide's intermediate-checkpoint comparison have no released method
outputs. We list these rather than re-implement them.

\section{Correcting a pair-resampled interval}
\label{sec:correction}

A practitioner who has already bootstrapped over pairs can correct the
interval without re-running anything. The pair-resampled variance of a pair
mean over $n$ candidates is $\zeta_2/\binom{n}{2}$, and the candidate-level  %% nolint: mathematical constant
variance is $2\bigl(2(n-2)\zeta_1+\zeta_2\bigr)/\bigl(n(n-1)\bigr)$. The  %% nolint: mathematical constant
standard-error ratio is therefore
\[
R(n)=\sqrt{1+2(n-2)\,\zeta_1/\zeta_2}.  %% nolint: mathematical constant
\]
It grows without bound in $n$: more candidates make pair resampling
\emph{more} anti-conservative, not less, because each candidate enters more
pairs. On the C4 comparison $R(\nCandidates)=\curveCFour$. Across the
\curveComparisons\ comparisons in this paper and the published re-tests, the
median $\zeta_1/\zeta_2$ is \curveZetaMedian, which gives
$R=\curveMedTen$, \curveMedTwentyFive\ and \curveMedHundred\ at ten,
twenty-five and a hundred candidates. The 10th--90th percentile range at  %% nolint: names which percentiles, values are macros
twenty-five candidates is \curveLowTwentyFive--\curveHighTwentyFive.  %% nolint: names which percentiles, values are macros
Subsampling the real candidates reproduces $R(m)$ to within
\curveCheckWorstLarge\% at twelve or more candidates. At \curveCheckSmallestM\
it misses by \curveCheckSmallest\%, a pre-registered prediction that failed:
with so few pairs the pair-resampled standard error is itself unstable. Where
$\zeta_1$ and $\zeta_2$ cannot be estimated, a pair-resampled interval at
twenty-five candidates should be widened by about a factor of two.

\section{What survives the correct test}
\label{sec:survives}

The paper makes \restatedPrimaryM\ inferential claims about selection methods:
comparisons of estimators and allocations against single-scale ranking on two
metrics, the lever-arm association, and the re-tests of published comparisons.
Each is restated under the calibrated random-candidate procedure, and the
family is corrected with Holm and with Benjamini--Yekutieli. BY is the right
one here, because every claim rests on the same \nCandidates\ recipes.
Appendix~\ref{app:claims} lists them all. Of the \restatedPrimaryM,
\restatedPrimaryNominal\ are nominally significant and \restatedPrimaryBY\
survives correction.

\emph{The survivor is the lever-arm dependence.} Excess projection error rises
with the logarithm of the lever arm: Spearman \doseRhoObserved\ against
observed targets and \doseRhoExpected\ under expected-error scoring, over
\doseDesigns\ designs. Those designs share recipes, so the design-level
$p$-value is not a valid test. A delete-one-recipe jackknife of the correlation
gives [\doseJackLow, \doseJackHigh] ($p=\doseJackP$).

\emph{What does not survive.} The ensemble's inferiority to single-scale ranking
is nominally significant for recipes in general but does not survive the
family correction. For these \nCandidates\ recipes it is the one ranker
comparison that survives correction in the fixed-candidate family
(\restatedFixedBY\ of \restatedFixedM; $p=\ensFixedP$). The optimised
allocation's excess at small budget fractions, one DataDecide scaling-law
variant, and Correct Prob at the two smallest scales are nominal only. No
comparison favours any estimator or allocation over single-scale ranking, even
nominally.

Certification is unaffected. The abstention rule certifies \certNormalN\ orderings
with no observed errors; under expected-error scoring those carry
\certNormalExpected\ errors in expectation, concentrated on the one certified pair
the reference itself resolves least well (probability \certNormalMinP). With the
Student-t critical value the rule certifies \certStudentTN\ and the expected
count is \certStudentTExpected.

\section{Fit structure, independent of the reference}
\label{sec:fit}

A separate set of findings concerns the fit itself: misspecification (residual
scatter \residRatio\ times the seed standard error), whether the misspecification
amplitude is shared across candidates, the insensitivity of rankings to a shared
monotone distortion, floor identifiability, the information content of
correlated checkpoints, resolvability, and disagreement between evaluation
metrics. None of these uses target-scale reference labels, so none is touched
by uncertainty propagation, and each still meets the criterion its source
states. Earlier drafts offered several of them as causes of the decision gap.
That gap is not significant once candidates are resampled, so we make no
causal claim; they are statements about the fit.

The floor result needed a correction of its own. The earlier profile fitted
$(A,\alpha)$ in log space at each fixed floor, so its residual was not
minimised on the response scale, and its admissibility band used run-level
rather than cell-mean variance. Refitting properly and taking the F-based
profile interval (critical value \floorFCrit), the C4 floor excludes zero for
\floorCFourExcl\ of \floorN\ candidates and the OLMES-error floor for
\floorErrExcl\ of \floorN, where the interval spans \floorErrWidthRel\% of the
feasible range. The information matrix is \floorCondRatio\ times worse
conditioned on the accuracy metric. The corrected C4 interval is wider than the
earlier one, which complicates the two-regime reading (\S\ref{sec:regimes}).

\section{Resolvability and the reference distribution}
\label{sec:resolvability}

A three-seed comparison does not have \maxDfAtThreeSeeds\ degrees of freedom except when the two
cells have equal variance; \maxDfAtThreeSeeds\ is the maximum. With per-pair Welch degrees of
freedom, the normal quantile under Bonferroni correction is anti-conservative by
a median factor of \tRatioMedian\ (10th--90th percentile \tRatioPTen--%  %% nolint: names which percentiles, values are macros
\tRatioPNinety), against \tRatioAssumed\ under the assumed
\maxDfAtThreeSeeds\ df. Reporting
significance on a three-seed grid with a Gaussian threshold is worse than
previously stated.

\section{A reproducibility failure in our own results}
\label{sec:provenance}

An earlier version of this work reported \shippedProjFlips\ projection ranking
errors and a corresponding single-scale count on the primary design. Those
counts do not regenerate. At the commit that wrote them, using that commit's own
code and byte-identical input data, the same detectors give
\recomputedProjFlips\ and \recomputedSingleFlips; the detector module is
unchanged since, and the input hashes match. The shipped file was produced by
uncommitted working state. We report the reproducible counts and state the
failure rather than quietly substituting them.

\section{Related work}
\label{sec:related}

\emph{Signal and Noise} (arXiv:2508.13144) defines decision accuracy as pairwise
agreement with a large-scale reference and shows a benchmark's signal-to-noise
ratio predicts it. They measure reference noise --- reporting a measured
accuracy range across the final checkpoints of its reference run --- but their
decision-accuracy estimator treats the reference ranking as ground truth, with
no filtering or weighting of pairs by whether the reference resolves them.
Reference uncertainty there predicts decision accuracy; here it is propagated
into it.

Northcutt et al.\ (arXiv:2103.14749) show corrected test-set labels can invert
model rankings. The structure is the same and the source of reference error is
not: theirs is systematic mislabelling, correctable by relabelling; ours is
sampling uncertainty in a finite-seed reference, which cannot be corrected at
any labelling effort, only quantified.

\section{Limitations}
\label{sec:limitations}

Evidence is one model suite and one intervention class (data recipes). Both
scoring rules assume unbiased target cell means. The determined-subset rule
inherits whatever multiplicity convention is chosen, and we report the whole
threshold curve rather than one cut precisely because that choice matters. The
provenance failure in \S\ref{sec:provenance} is the eighth instance of a
measurement artifact in this project; the practice record and its detection
paths are maintained alongside the code.

The power calculation holds the variance components at their estimates from
one candidate pool and uses a normal approximation. It says what a comparison
of this kind needs, not what every suite needs. The candidate-level estimator
we report is a declared deviation from the pre-registered one, made after
simulation showed the pre-registered bootstrap to be conservative here; both are
reported. The allocation study's gradient solver stops \thmSolverExcess\% above
the true optimum of the design problem on this ladder
(Appendix~\ref{app:design}). We did not re-run the allocation study, because a
variance difference of that size is far below the margins it reports. The
scaling-law re-test is exploratory while the discrepancy between the released
decision accuracies and our recomputation from the released predictions is
under investigation.

The calibration benchmark is only as good as its assumptions. Its truth is the
observed seed-mean trajectories and its noise is Gaussian. Its random-candidate
worlds come from a smoothed bootstrap over only \nCandidates\ recipes, and its
critical values are estimated from \calFixedWorlds\ and \calRandomWorlds\
worlds. They were calibrated on C4. For OLMES error and the published re-tests
they are transported, not re-estimated. EB shrinkage enters the benchmark with
its strength frozen at the real-data estimate, which omits the variability of
estimating it. Pooled calibration leaves per-comparator held-out coverage
between \calFixedCovMin\ and \calJackCovMax, not uniformly at nominal. Every
p-value in Appendix~\ref{app:claims} inherits these assumptions.

Resolvability figures for suites without seed replicates transfer a seed
standard deviation from another scale. Whether seed noise changes with model
size is untested, and cannot be tested on current public infrastructure. A
trend needs replicates at three or more sizes on one recipe. Marin's Delphi
ladder, the newest public suite, has them only at 1e21 and 1e22 FLOPs (three  %% nolint: compute-budget labels and a seed count read from the model listing, not measured values
seeds each); every other budget, including its largest, has a single run.

\appendix
\section{Design for estimation and design for decision coincide}
\label{app:design}

\paragraph{Setting.} Candidates $i$ are fitted on one shared design: budgets
$C_1,\dots,C_k$ with $n_j$ runs each. Write $x_j=(1,-\log C_j)^\top$ and  %% nolint: mathematical constant
$A(n)=\sum_j n_j x_j x_j^\top$, assumed nonsingular. Observations are
$y_{ijr}=x_j^\top\theta_i+\varepsilon_{ijr}$ in log space, with $\theta_i$
estimated by least squares.

\paragraph{Assumptions.} (A1) Every candidate uses the same design $n$.
(A2) The errors have mean zero and variance $\sigma^2$, are independent across  %% nolint: mathematical constant
runs, and are independent across candidates. (A3) The model is well specified:
there is no deviation $b(C)$ from the linear form.

\paragraph{Theorem.} Under (A1)--(A3), for every pair $(i,i')$ and every design $n$,
\[
\operatorname{Var}\bigl(x_*^\top(\hat\theta_i-\hat\theta_{i'})\bigr)
 = 2\operatorname{Var}\bigl(x_*^\top\hat\theta_i\bigr)  %% nolint: mathematical constant
 = 2\sigma^2 v(n),\qquad v(n)=x_*^\top A(n)^{-1}x_*,  %% nolint: mathematical constant
\]
so over any feasible set of designs the variance of the target gap and the
variance of the target level have the same minimisers, for all pairs at once. If
the errors are also Gaussian, the probability that the estimated gap has the
wrong sign is $\Phi\bigl(-|\Delta_{ii'}|/\sqrt{2\sigma^2 v(n)}\bigr)$, which is  %% nolint: mathematical constant
increasing in $v(n)$. The level-optimal design therefore minimises every pair's
mis-selection probability simultaneously.

\paragraph{Proof.} Least squares gives
$\hat\theta_i-\theta_i=A(n)^{-1}\sum_j x_j\sum_r\varepsilon_{ijr}$, so by (A2)  %% nolint: mathematical constant
\[  %% nolint: mathematical constant
\operatorname{Cov}(\hat\theta_i)=\sigma^2A(n)^{-1}\bigl(\sum_j n_jx_jx_j^\top\bigr)A(n)^{-1}=\sigma^2A(n)^{-1}.  %% nolint: mathematical constant
\]
Independence across candidates gives
$\operatorname{Cov}(\hat\theta_i-\hat\theta_{i'})=2\sigma^2A(n)^{-1}$. The pair  %% nolint: mathematical constant
variance is therefore twice the level variance, a factor that does not depend on
$n$, so the minimisers coincide. By (A3) the estimated gap is unbiased for
$\Delta_{ii'}$. Under Gaussian errors it is distributed
$\mathcal N(\Delta_{ii'},2\sigma^2v(n))$, and the wrong-sign probability is as  %% nolint: mathematical constant
stated. $\Phi$ is increasing and $|\Delta_{ii'}|$ does not depend on $n$, so
that probability is increasing in $v(n)$. \hfill$\square$

\paragraph{What the assumptions exclude.} Budget-dependent variance shared by
all candidates is harmless: it changes $A$ to a weighted matrix, the same for
every candidate. The conclusion fails when (A1) fails, because candidates on
different designs have different $A$. It fails when the correlation between
candidates' errors varies with budget, because the factor $2$ becomes a  %% nolint: mathematical constant
design-dependent matrix. A constant correlation $\rho$ only rescales the factor
to $2(1-\rho)$. It also fails when (A3) fails: a deviation $b(C)$ adds the bias  %% nolint: mathematical constant
$x_*^\top A(n)^{-1}X^\top W b$, which depends on $n$. That last case is the one  %% nolint: mathematical constant
the measurements in \S\ref{sec:fit} show holds in practice. The theorem is about
design; it says nothing about whether projection should beat single-scale
ranking.

\paragraph{The region objective.} A target-region criterion that averages the
prediction variance over a region around $C_*$ is not covered by the theorem.
Numerically it is nearly equivalent. On the DataDecide ladder, its correlation
with the point objective over \thmDesigns\ random designs is \thmCorr, and the
design optimal for a region spanning three decades costs \thmExcessThree\% extra
point variance at $C_*$. These optima come from an exhaustive search over two-
and three-point supports. A gradient solver in run-count space stops
\thmSolverExcess\% short of the optimum on the same ladder, because run counts
span orders of magnitude across the budgets.


\section{Every inferential claim, restated}
\label{app:claims}

\begin{table}[h]
\centering
\scriptsize
\begin{tabular}{p{7.2cm}rrrl}  %% nolint: column width
\toprule
claim (random-candidate estimand) & estimate & calibrated interval & $p$ & verdict \\
\midrule
\restatedPrimaryRows
\bottomrule
\end{tabular}
\caption{The paper's inferential claims under the calibrated random-candidate
procedure. Verdict: Holm or BY if the claim survives that correction across the
family, nominal if only $p<0.05$, otherwise a dash.}  %% nolint: the stated significance level
\label{tab:claims}
\end{table}

\begin{table}[h]
\centering
\scriptsize
\begin{tabular}{p{7.2cm}rrrl}  %% nolint: column width
\toprule
claim (fixed-candidate estimand) & estimate & calibrated interval & $p$ & verdict \\
\midrule
\restatedFixedRows
\bottomrule
\end{tabular}
\caption{Ranker comparisons under the fixed-candidate estimand, scoped to these
recipes, corrected within their own family.}
\label{tab:claims-fixed}
\end{table}
\end{document}
